% संपूर्ण मराठी ग्रंथातील प्रकरण 7: विधानीय तर्कशास्त्र: विन्यासमीमांसा आणि चिन्हार्थमीमांसा
% हा संपादनयोग्य स्रोतखंड आहे. संकलनासाठी संपूर्ण स्रोतसंग्रहातील
% अक्षररूपे, पूर्वभाग, चिन्हव्याख्या आणि आकृती-स्रोत आवश्यक आहेत.
% मूळ: Open Logic Project; मजकूर आणि रूपांतर CC BY 4.0.
% यंत्रानुवाद, दुरुस्ती आणि तपासणी: OpenAI Codex —
% GPT-5.6 Sol आणि GPT-6 Sol, दोन्ही Ultra effort.
% दुरुस्ती-पुनरावलोकन: GPT-6.1 Sol, Ultra effort.
\chapter{विधानीय तर्कशास्त्र: विन्यासमीमांसा आणि चिन्हार्थमीमांसा}\label{pl:syn::chap}
\begin{editorial}
  या भागात अभिजात विधानीय तर्कशास्त्रावरील सामग्री आहे. पहिले
  प्रकरण तुलनेने प्राथमिक स्वरूपाचे आहे आणि त्यात केवळ व्याख्या
  व निष्पत्त्या दिल्या आहेत; अनेक सिद्धता पूर्ण केलेल्या नसून त्या
  सरावासाठी सोडल्या आहेत. सिद्धता-पद्धतींवरील आणि संपूर्णता
  प्रमेयावरील सामग्री प्रथम-क्रम तर्कशास्त्राच्या भागातून, ``FOL''
  हा टॅग false असा ठेवून, समाविष्ट केली आहे. त्यामुळे विधेये,
  पदे आणि संख्यापक यांच्याशी संबंधित सर्व मजकूर वगळला जातो;
  तसेच संरचना~$\Struct{M}$ या संरचनांविषयीच्या चर्चेऐवजी
  सत्य-मूल्यांकने~$\pAssign{v}$ या सत्य-मूल्यांकनांविषयी चर्चा येते.

  या भागात अधिक तपशील समाविष्ट करण्यासाठी त्याचा विस्तार करणे,
  आणि सत्यता-फलनात्मक संपूर्णतेसारखे पुढील विषय व निष्पत्त्या जोडणे,
  असे नियोजन आहे.
\end{editorial}
\begin{editorial}
  हा केवळ व्याख्यांचा अतिशय संक्षिप्त आढावा आहे. सूत्रांवरील
  विगमनाने सिद्धता करण्याचा सुलभ परिचय आणि आणखी बरीच उदाहरणे
  देऊन त्याचा विस्तार करायला हवा.
\end{editorial}






\section{परिचय}\label{pl:syn:int:sec}

विधानीय तर्कशास्त्रात विधानीय संयोजके
$\lnot$, $\land$, $\lor$, $\lif$ आणि $\liff$ वापरून
विधानीय चलांपासून बनवलेल्या सूत्रांचा विचार
केला जातो. अंतःप्रेरणेने पाहता, विधानीय चल~$p$
हे सत्य किंवा असत्य असलेल्या एखाद्या वाक्याच्या किंवा विधानाच्या
जागी येते. सूत्र मधील विधानीय चलाचे
``सत्यतामूल्य'' ठरले की, विधानीय संयोजके वापरून त्यांच्यापासून
बनवलेल्या कोणत्याही सूत्रांचे सत्यतामूल्यही ठरते. विधानीय
तर्कशास्त्राला आपण \emph{सत्यता-फलनात्मक} म्हणतो, कारण त्याची
चिन्हार्थमीमांसा सत्यतामूल्यांच्या फलनांनी दिली जाते. विशेषतः,
विधानीय तर्कशास्त्रात सत्य आणि असत्य यांचे पुढील प्रकारचे निर्धारण
आपण विचारात घेत नाही: उदाहरणार्थ, एखादी गोष्ट केवळ परिस्थितीवश
सत्य असण्याऐवजी अनिवार्यपणे सत्य आहे का, ती सत्य आहे हे ज्ञात आहे
का, किंवा ती आत्ता सत्य आहे की पूर्वी सत्य होती अथवा पुढे सत्य
असेल. आपण फक्त सत्य ($\True$) आणि असत्य ($\False$) ही दोन
सत्यतामूल्ये विचारात घेतो. म्हणून एखादे विधान सत्यही नसते व
असत्यही नसते, किंवा ते केवळ अर्धसत्य असते, ही शक्यता चर्चेतून
वगळतो. तसेच, ज्यांच्यापासून बनवलेल्या सूत्राचे सत्यतामूल्य
त्याच्या भागांच्या सत्यतामूल्यांनी पूर्णपणे ठरते (त्याच्या अर्थाने,
उदाहरणार्थ, ठरत नाही), अशाच संयोजकांवर आपण लक्ष केंद्रित करतो.
विशेषतः, इंग्रजीतील सशर्त विधानांचे सत्यतामूल्य या अर्थाने
सत्यता-फलनात्मक असते की नाही, हा वादग्रस्त प्रश्न आहे. वास्तविक
अभिव्यंजन~$\lif$ सत्यता-फलनात्मक आहे; सत्यता-फलनात्मक नसलेल्या सशर्त
विधानांचा विचार इतर तर्कपद्धती करतात.

सत्यता-फलनात्मक विधानीय तर्कशास्त्राची उपपत्ती आणि अधिउपपत्ती
विकसित करण्यासाठी, आपण प्रथम त्याच्या व्यंजकांची विन्यासमीमांसा
आणि चिन्हार्थमीमांसा परिभाषित केली पाहिजे. संयोजके वापरून
विधानीय चलांपासून सूत्रे तयार करण्याची एक
पद्धत आपण वर्णन करू. पर्यायी व्याख्या शक्य आहेत. इतर पद्धती वेगळी
चिन्हे निवडतील, संयोजकांचे वेगळे संच मूलभूत म्हणून निवडतील आणि
कंसांचा वेगळ्या रीतीने वापर करतील (किंवा तथाकथित पोलिश
संकेतपद्धतीप्रमाणे कंस अजिबात वापरणार नाहीत). तरी सर्व दृष्टिकोनांत
सामाईक बाब ही की रचनानियम सूत्रांचा संच \emph{विगमनाने}
परिभाषित करतात. हे योग्य रीतीने केल्यास, रचनानियमांनुसार प्रत्येक
व्यंजक मूलतः एकाच रीतीने तयार होऊ शकते. त्यामुळे तयार होणारी
व्यंजके \emph{एकमेव रीतीने वाचनीय} असतात. परिणामी, त्यांची ही
विगमनाधारित व्याख्या आपल्याला त्याच पद्धतीने---विगमनाधारित
व्याख्येने---त्या व्यंजकांना अर्थ देता येतो, याची खात्री देते.

व्यंजकांना अर्थ देणे हा चिन्हार्थमीमांसेचा विषय आहे. विधानीय
तर्कशास्त्राच्या चिन्हार्थमीमांसेतील मध्यवर्ती संकल्पना म्हणजे
सत्य-मूल्यांकन मधील पूर्ति. सत्य-मूल्यांकन~$\pAssign{v}$ हे
विधानीय चलांना $\True$, $\False$ ही सत्यतामूल्ये
नेमून देते. कोणतेही सत्य-मूल्यांकन कोणत्याही सूत्र~$\varphi$
साठी $\pValue{v}(\varphi)$ हे सत्यतामूल्य ठरवते. सूत्राची
सत्य-मूल्यांकन~$\pAssign{v}$ मध्ये पूर्ति तेव्हा आणि केवळ तेव्हाच
होते जेव्हा $\pValue{v}(\varphi) = \True$; हे आपण $\pSat{v}{\varphi}$ असे
लिहितो. तार्किक संयोजकांची सत्यता-फलने वापरून, उदाहरणार्थ
$\varphi \land \psi$ ची पूर्ति $\varphi$ आणि~$\psi$ यांची पूर्ति होते (किंवा
होत नाही) यावरून परिभाषित करून, हा संबंधही~$\varphi$ च्या रचनेवरील
विगमनाने परिभाषित करता येतो.

वाक्यांसाठी असलेल्या $\pSat{v}{\varphi}$ या पूर्ति-संबंधाच्या आधारे
आपण सर्वतः सत्यता, तार्किक निष्पन्नता आणि पूर्ततायोग्यता या मूलभूत
चिन्हार्थविषयक संकल्पना परिभाषित करू शकतो. सूत्र हे सर्वतः
सत्य सूत्र, $\Entails \varphi$, असते, जर प्रत्येक सत्य-मूल्यांकन त्याची
पूर्ति करत असेल, म्हणजे कोणत्याही~$\pAssign{v}$ साठी
$\pValue{v}(\varphi) = \True$ असेल. सूत्रांच्या एखाद्या संचापासून,
$\Gamma \Entails \varphi$, ते तार्किकरीत्या निष्पन्न होते, जर~$\Gamma$
मधील सर्व सूत्रांची पूर्ति करणारे प्रत्येक सत्य-मूल्यांकन
~$\varphi$ चीही पूर्ति करत असेल. तसेच, सूत्रांचा एखादा संच
पूर्ततायोग्य असतो, जर एखादे सत्य-मूल्यांकन त्यातील सर्व सूत्रांची एकाच वेळी पूर्ति करत असेल. सूत्रे विगमनाने
परिभाषित केलेली असल्यामुळे, आणि पूर्तिही सूत्रांच्या रचनेवरील
विगमनाने परिभाषित केलेली असल्यामुळे, आपल्या चिन्हार्थमीमांसेचे
गुणधर्म सिद्ध करण्यासाठी आणि परिभाषित केलेल्या चिन्हार्थविषयक
संकल्पनांचा परस्परसंबंध दाखवण्यासाठी आपण विगमन वापरू शकतो.









\section{विधानीय सूत्र}\label{pl:syn:fml:sec}

विधानीय तर्कशास्त्रातील सूत्रे ही
\emph{विधानीय चले}, विधानीय
  स्थिर~$\lfalse$
    आणि विधानीय स्थिर~$\ltrue$ यांपासून \emph{तार्किक
  संयोजके} वापरून तयार होतात.

\begin{enumerate}
\item प्रगणनीय संच~$\PVar$: त्यातील विधानीय चले
  $\Obj p_0$, $\Obj p_1$, \dots
\item असत्यता साठीचे विधानीय स्थिर~$\lfalse$.
\item सत्यता साठीचे विधानीय स्थिर~$\ltrue$.
\item तार्किक संयोजके:
  \startycommalist
  \ycomma $\lnot$ (नकरण)%
  \ycomma $\land$ (संधियोग)%
  \ycomma $\lor$ (विकल्पयोग)%
  \ycomma $\lif$ (शर्त)%
  \ycomma $\liff$ (द्विशर्त)%
\item विरामचिन्हे: (, ), आणि स्वल्पविराम.
\end{enumerate}

विधानीय तर्कशास्त्राची ही भाषा आपण $\Lang L_0$ ने दर्शवतो.

.




\begin{intro}
वर आपण वापरलेल्या संज्ञांपेक्षा आणि चिन्हांपेक्षा वेगळ्या संज्ञा व
चिन्हे तुम्हाला परिचित असू शकतात. तर्कशास्त्राच्या ग्रंथांत (आणि
अध्यापनात) ``नकरण'' यासाठी सामान्यतः $\sim$, $\neg$, आणि~!\ यांपैकी
एखादे, तर ``संधियोग'' यासाठी $\wedge$, $\cdot$, आणि $\&$ यांपैकी
एखादे चिन्ह वापरतात. ``सशर्त (संकेतार्थक)'' किंवा ``अभिव्यंजन'' यांसाठी
$\rightarrow$, $\Rightarrow$, आणि $\supset$ ही चिन्हे सामान्यतः
वापरली जातात.
``द्विसशर्त'', ``द्वि-अभिव्यंजन'' किंवा
  ``(वास्तविक) सममूल्यता'' यांसाठी $\leftrightarrow$,
  $\Leftrightarrow$, आणि $\equiv$ ही चिन्हे वापरतात.
$\lfalse$ या चिन्हाला वेगवेगळ्या ठिकाणी
  ``असत्यता'', ``फाल्सम'', ``विसंगती'' किंवा ``तळ'' म्हणतात.
$\ltrue$ या चिन्हाला वेगवेगळ्या ठिकाणी
  ``सत्यता'', ``वेरम'' किंवा ``शिखर'' म्हणतात.
\end{intro}

\begin{defn}[सूत्र]
\label{pl:syn:fml:defn:formulas}
विधानीय तर्कशास्त्रातील \emph{सूत्रे} चा संच~$\Frm[L_0]$
पुढीलप्रमाणे विगमनाने परिभाषित केला जातो:
\begin{enumerate}
\item $\lfalse$ हे आण्विक सूत्र आहे.

\item $\ltrue$ हे आण्विक सूत्र आहे.

\item प्रत्येक विधानीय चल~$\Obj p_i$ हे आण्विक
  सूत्र आहे.

\item $\varphi$ हे सूत्र असेल, तर $\lnot \varphi$ हे
  सूत्र आहे.

\item $\varphi$ आणि $\psi$ ही सूत्रे असतील, तर $(\varphi \land
  \psi)$ हे सूत्र आहे.

\item $\varphi$ आणि $\psi$ ही सूत्रे असतील, तर $(\varphi \lor \psi)$
  हे सूत्र आहे.

\item $\varphi$ आणि $\psi$ ही सूत्रे असतील, तर $(\varphi \lif \psi)$
  हे सूत्र आहे.

\item $\varphi$ आणि $\psi$ ही सूत्रे असतील, तर $(\varphi \liff \psi)$
  हे सूत्र आहे.

\item इतर काहीही सूत्र नाही.
\end{enumerate}
\end{defn}

\begin{explain}
सूत्रांची ही व्याख्या \emph{विगमनाधारित व्याख्या} आहे. मूलतः,
आपण सूत्रांचा संच अनंतपणे अनेक टप्प्यांत तयार करतो. आरंभीच्या
टप्प्यात आपण सर्व आण्विक सूत्रांना सूत्रे घोषित करतो; हे व्याख्येतील
पहिल्या काही बाबींशी, म्हणजे
$\ltrue$, %
$\lfalse$, %
$\Obj p_i$ यांच्या बाबींशी जुळते. म्हणून ``आण्विक सूत्र''
म्हणजे या रूपातील कोणतेही सूत्र.

व्याख्येतील इतर बाबी आधीच तयार केलेल्या सूत्रांपासून नवी
सूत्रे तयार करण्याचे नियम देतात. दुसऱ्या टप्प्यात या नियमांनी
आण्विक सूत्रांपासून सूत्रे तयार करता येतात. तिसऱ्या
टप्प्यात आण्विक सूत्रे आणि दुसऱ्या टप्प्यात मिळालेली सूत्रे यांपासून
नवी सूत्रे तयार करतो, आणि ही प्रक्रिया पुढे चालू राहते. अशा एखाद्या
टप्प्यात शेवटी तयार होणारी प्रत्येक गोष्ट सूत्र असते; इतर
काहीही सूत्र नसते.
\end{explain}

द्विस्थानी संयोजक~$\ast$ वापरून $\psi$, $\chi$ पासून तयार केलेले
सूत्र $(\psi \ast \chi)$ लिहिताना आपण अनेकदा सर्वांत बाहेरची कंसांची
जोडी वगळून केवळ~$\psi \ast \chi$ असे लिहू.



\begin{defn}[विन्यासात्मक अनन्यता]
$\ident$ हे चिन्ह चिन्हमालांमधील विन्यासात्मक अनन्यता व्यक्त करते;
म्हणजे $\varphi \ident \psi$ तेव्हा आणि केवळ तेव्हाच, जेव्हा $\varphi$ आणि $\psi$
या समान लांबीच्या चिन्हमाला असतात आणि प्रत्येक स्थानी त्यांच्यात तेच
चिन्ह असते.
\end{defn}

$\ident$ या चिन्हाच्या दोन्ही बाजूंना जोडणीने मिळालेल्या चिन्हमाला
येऊ शकतात. उदाहरणार्थ, $\varphi \ident (\psi \lor \chi)$ याचा अर्थ असा:
आरंभीचा कंस, चिन्हमाला~$\psi$, $\lor$ हे चिन्ह, चिन्हमाला~$\chi$ आणि
शेवटचा कंस या क्रमाने जोडून मिळणारी चिन्हमाला आणि चिन्हमाला~$\varphi$
या एकच आहेत. असे असेल, तर $\varphi$ चे पहिले चिन्ह आरंभीचा कंस आहे;
$\varphi$ मध्ये दुसऱ्या चिन्हापासून सुरू होणारी $\psi$ ही उपचिन्हमाला आहे;
तिच्यानंतर $\lor$ येते; इत्यादी बाबी आपल्याला माहीत होतात.









\section{पूर्वतयारी}\label{pl:syn:pre:sec}

\begin{thm}[\emph{सूत्रांवरील विगमनाचे तत्त्व}]
  \label{pl:syn:pre:thm:induction}
  एखादा गुणधर्म~$P$ सर्व आण्विक सूत्रांसाठी लागू होत असेल आणि
  पुढील अटी पूर्ण करत असेल:
  \begin{enumerate}
    \item तो~$\varphi$ साठी लागू असताना $\lnot \varphi$ साठीही
      लागू होतो;
    \item तो~$\varphi$ आणि~$\psi$ यांच्यासाठी लागू असताना
      $(\varphi \land \psi)$ साठीही लागू होतो;
    \item तो~$\varphi$ आणि~$\psi$ यांच्यासाठी लागू असताना
      $(\varphi \lor \psi)$ साठीही लागू होतो;
    \item तो~$\varphi$ आणि~$\psi$ यांच्यासाठी लागू असताना
      $(\varphi \lif \psi)$ साठीही लागू होतो;
    \item तो~$\varphi$ आणि~$\psi$ यांच्यासाठी लागू असताना
      $(\varphi \liff \psi)$ साठीही लागू होतो;
  \end{enumerate}
  तर $P$ सर्व सूत्रांसाठी लागू होतो.
\end{thm}

\begin{proof}
  गुणधर्म~$P$ असलेल्या सर्व सूत्रांचा संग्रह~$S$ असू द्या.
  स्पष्टपणे $S \subseteq \Frm[L_0]$. \ref{pl:syn:fml:defn:formulas} मधील
  सर्व अटी $S$~पूर्ण करतो: त्यात सर्व आण्विक सूत्रे आहेत आणि
  तो संकारकांखाली बंद आहे.  $\Frm[L_0]$ हा असा सर्वांत लहान
  वर्ग आहे; म्हणून $\Frm[L_0] \subseteq S$. त्यामुळे $\Frm[L_0] = S$,
  आणि प्रत्येक सूत्राला गुणधर्म~$P$ आहे.
\end{proof}

\begin{prop}\label{pl:syn:pre:prop:balanced}
   $\Frm[L_0]$ मधील प्रत्येक सूत्र \emph{संतुलित} आहे; म्हणजे
   त्यात जेवढे डावे कंस आहेत तेवढेच उजवे कंस आहेत.
\end{prop}

\begin{prob}
\ref{pl:syn:pre:prop:balanced} सिद्ध करा.
\end{prob}

\begin{prop} \label{pl:syn:pre:prop:noinit}
कोणत्याही सूत्राचा कोणताही उचित आरंभीचा खंड सूत्र नसतो.
\end{prop}

\begin{prob}
  \ref{pl:syn:pre:prop:noinit} सिद्ध करा.
\end{prob}

\begin{prop}[एकमेव वाचनीयता]
$ \Frm[L_0]$ मधील कोणत्याही सूत्र~$\varphi$ चे पुढीलपैकी नेमक्या एका
रूपात रचनाविश्लेषण करता येते:
\begin{enumerate}
\item $\lfalse$.

\item $\ltrue$.

\item एखाद्या $\Obj p_n \in  \PVar$ साठी $\Obj p_n$.

\item एखाद्या सूत्र~$\psi$ साठी $\lnot \psi$.

\item काही सूत्रे $\psi$ आणि~$\chi$ यांच्यासाठी
  $(\psi \land \chi)$.

\item काही सूत्रे $\psi$ आणि~$\chi$ यांच्यासाठी
  $(\psi \lor \chi)$.

\item काही सूत्रे $\psi$ आणि~$\chi$ यांच्यासाठी
  $(\psi \lif \chi)$.

\item काही सूत्रे $\psi$ आणि~$\chi$ यांच्यासाठी
  $(\psi \liff \chi)$.
\end{enumerate}
शिवाय, हे रचनाविश्लेषण \emph{एकमेव} असते.
\end{prop}

\begin{proof}
$\varphi$ वरील विगमनाने सिद्ध करू. उदाहरणार्थ, $\varphi$ चे $(\psi \lif \chi)$ आणि
$(\psi' \lif \chi')$ अशी दोन भिन्न रचनाविश्लेषणे आहेत असे समजा. मग $\psi$
आणि $\psi'$ एकच असले पाहिजेत (अन्यथा त्यांपैकी एक हा दुसऱ्याचा उचित
आरंभीचा खंड असेल); म्हणून $\varphi$ ची ही दोन रचनाविश्लेषणे भिन्न असतील,
तर त्याचे कारण असेच असले पाहिजे की $\chi$ आणि $\chi'$ ही एकाच
चिन्हमालेची भिन्न रचनाविश्लेषणे आहेत; पण विगमन गृहीतकानुसार ते अशक्य
आहे.
\end{proof}


\begin{defn}[एकरूप आदेशन]
$\varphi$ आणि $\psi$ ही सूत्रे असतील आणि $\Obj p_i$ हे विधानीय चल असेल, तर $\varphi$ मधील $\Obj p_i$ ची प्रत्येक
आस्थिति $\psi$ च्या एका आस्थितीने बदलल्यावर मिळणारा परिणाम
$\Subst{\varphi}{\psi}{\Obj p_i}$ ने दर्शवला जातो; त्याचप्रमाणे,
$\Obj p_1$, \dots,~$\Obj p_n$ यांची सूत्रे $\psi_1$,
\dots,~$\psi_n$ यांनी एकाच वेळी केलेले आदेशन
$\SSubst{\varphi}{\subst{\psi_1}{\Obj p_1},\dots,\subst{\psi_n}{\Obj p_n}}$
ने दर्शवले जाते.
\end{defn}

\begin{prob} खालील पाच सूत्रांपैकी प्रत्येकासाठी, ते
 सूत्र \( \Subst{\varphi}{\psi}{\Obj p_i} \) या आदेशनाच्या रूपात व्यक्त
 करता येते का ते ठरवा; येथे \( \varphi \) हे अनुक्रमे (i) \( \Obj p_0 \);
 (ii) \( ( \lnot \Obj p_0 \land \Obj p_1) \); आणि (iii)
 \( ( ( \lnot \Obj p_0 \lif \Obj p_1 ) \land \Obj p_2 ) \) आहे.
 प्रत्येक बाबतीत संबद्ध आदेशन स्पष्ट करा.
  \begin{enumerate}
    \item \( \Obj p_1 \)
    \item \( ( \lnot \Obj p_0 \land \Obj p_0 ) \)
    \item \( ( ( \Obj p_0 \lor \Obj p_1 ) \land \Obj p_2 )  \)
    \item \( \lnot ( ( \Obj p_0 \lif \Obj p_1 ) \land \Obj p_2 ) \)
    \item \( (( \lnot ( \Obj p_0 \lif \Obj p_1 ) \lif ( \Obj p_0 \lor \Obj p_1 )) \land \lnot ( \Obj p_0 \land \Obj p_1 )) \)
  \end{enumerate}
\end{prob}

\begin{prob}
$\Subst{\varphi}{\psi}{p}$ ची विगमनाधारित, गणितीयदृष्ट्या काटेकोर व्याख्या
द्या.
\end{prob}









\section{रचनाक्रमिका}\label{pl:syn:fseq:sec}

सूत्रांची विगमनाधारित व्याख्या देणे आणि त्याला पूरक तंत्र म्हणून
सूत्रांचे गुणधर्म विगमनाने सिद्ध करणे हा सुबक व कार्यक्षम
दृष्टिकोन आहे. तथापि, सूत्रांच्या रचनेचा तळापासून वर जाणारा,
टप्प्याटप्प्याचा दृष्टिकोनही उपयुक्त ठरू शकतो; येथे आपण तो
\emph{रचनाक्रमिका} या संकल्पनेच्या साहाय्याने मांडतो.

\begin{defn}[सूत्रांसाठी रचनाक्रमिका]
\label{pl:syn:fseq:defn:fseq-frm}
भाषा~$\Lang L_0$ मधील चिन्हांपासून बनलेल्या चिन्हमालांची सांत क्रमिका
$\tuple{\varphi_0,\dotsc,\varphi_n}$ ही $\varphi$ साठी \emph{रचनाक्रमिका} असते, जर
$\varphi \ident \varphi_n$ आणि प्रत्येक $i \leq n$ साठी एकतर $\varphi_i$ हे आण्विक
सूत्र असेल किंवा पुढीलपैकी एखादी अट पूर्ण करणारे $j,k < i$ अस्तित्वात
असतील:
\begin{enumerate}
    \item $\varphi_i \ident \lnot \varphi_j$.%
    \item $\varphi_i \ident (\varphi_j \land \varphi_k)$.%
    \item $\varphi_i \ident (\varphi_j \lor \varphi_k)$.%
    \item $\varphi_i \ident (\varphi_j \lif \varphi_k)$.%
    \item $\varphi_i \ident (\varphi_j \liff \varphi_k)$.%
\end{enumerate}
\end{defn}

\begin{ex}
\[\tuple{
    \Obj p_0,
    \Obj p_1,
    (\Obj p_1 \land \Obj p_0),
    \lnot (\Obj p_1 \land \Obj p_0)
}
\]
ही $\lnot (\Obj p_1 \land \Obj p_0)$ ची एक रचनाक्रमिका आहे; तसेच
पुढील क्रमिकाही तिची रचनाक्रमिका आहे:
\[\tuple{
    \Obj p_0,
    \Obj p_1,
    \Obj p_0,
    (\Obj p_1 \land \Obj p_0),
    (\Obj p_0 \lif \Obj p_1),
    \lnot (\Obj p_1 \land \Obj p_0)
}.
\]
दुसऱ्या उदाहरणावरून दिसते की रचनाक्रमिकेत `अडगळ' असू शकते: अनावश्यक
असलेली किंवा रचनेत काहीही हातभार न लावणारी सूत्रे.
\end{ex}

\begin{prop}
\label{pl:syn:fseq:prop:formed}
$\Frm[L_0]$ मधील प्रत्येक सूत्र~$\varphi$ ला एक रचनाक्रमिका असते.
\end{prop}

\begin{proof}
$\varphi$ आण्विक आहे असे समजा. मग क्रमिका~$\tuple{\varphi}$ ही $\varphi$ साठी
रचनाक्रमिका आहे.
आता $\psi$ आणि~$\chi$ यांना अनुक्रमे $\tuple{\psi_0,\dotsc,\psi_n}$ आणि
$\tuple{\chi_0,\dotsc,\chi_m}$ या रचनाक्रमिका आहेत असे समजा.
\begin{enumerate}
  \item जर $\varphi \ident \lnot \psi$, तर
      $\tuple{\psi_0,\dotsc,\psi_n,\lnot \psi_n}$ ही $\varphi$ साठी
      रचनाक्रमिका आहे.
  \item जर $\varphi \ident (\psi \land \chi)$, तर
      $\tuple{\psi_0,\dotsc,\psi_n,\chi_0,\dotsc,\chi_m,(\psi_n \land \chi_m)}$
      ही $\varphi$ साठी रचनाक्रमिका आहे.
  \item जर $\varphi \ident (\psi \lor \chi)$, तर
      $\tuple{\psi_0,\dotsc,\psi_n,\chi_0,\dotsc,\chi_m,(\psi_n \lor \chi_m)}$
      ही $\varphi$ साठी रचनाक्रमिका आहे.
  \item जर $\varphi \ident (\psi \lif \chi)$, तर
      $\tuple{\psi_0,\dotsc,\psi_n,\chi_0,\dotsc,\chi_m,(\psi_n \lif \chi_m)}$
      ही $\varphi$ साठी रचनाक्रमिका आहे.
  \item जर $\varphi \ident (\psi \liff \chi)$, तर
      $\tuple{\psi_0,\dotsc,\psi_n,\chi_0,\dotsc,\chi_m,(\psi_n \liff \chi_m)}$
      ही $\varphi$ साठी रचनाक्रमिका आहे.
\end{enumerate}
सूत्रांवरील विगमनाच्या तत्त्वानुसार प्रत्येक सूत्राला एक
रचनाक्रमिका असते.
\end{proof}

याचे उलट विधानही आपण सिद्ध करू शकतो. ते महत्त्वाचे आहे, कारण त्यामुळे
सूत्रे परिभाषित करण्याच्या आपल्या दोन्ही पद्धती सममूल्य आहेत, म्हणजे
त्या समान परिणाम देतात, हे दिसते. तसेच रचनाक्रमिकांच्या लांबीवर साधे
विगमन वापरून सूत्रांविषयीची प्रमेये सिद्ध करता येतात.

\begin{lem}
\label{pl:syn:fseq:lem:fseq-init}
$\tuple{\varphi_0,\dotsc,\varphi_n}$ ही $\varphi_n$ साठी रचनाक्रमिका आहे आणि
$k \leq n$ असे समजा. मग $\tuple{\varphi_0,\dotsc,\varphi_k}$ ही $\varphi_k$ साठी
रचनाक्रमिका आहे.
\end{lem}

\begin{proof}
अभ्यासप्रश्न.
\end{proof}

\begin{thm}
\label{pl:syn:fseq:thm:fseq-frm-equiv}
$\Frm[L_0]$ हा भाषा~$\Lang L_0$ मधील ज्या सर्व चिन्हमालांना
रचनाक्रमिका आहे, त्यांचा संच आहे.
\end{thm}

\begin{proof}
भाषा~$\Lang L_0$ मधील ज्या सर्व चिन्हमालांना रचनाक्रमिका आहे, त्यांचा
संच~$F$ असू द्या. $\Frm[L_0] \subseteq F$ हे
\ref{pl:syn:fseq:prop:formed} मध्ये आपण पाहिले आहे; म्हणून आता उलट
समावेश सिद्ध करू.

$\varphi$ ला $\tuple{\varphi_0,\dotsc,\varphi_n}$ ही रचनाक्रमिका आहे असे समजा.
$n$ वरील प्रबल विगमनाने $\varphi \in \Frm[L_0]$ हे सिद्ध करू. लांबी
$m < n$ असलेली रचनाक्रमिका ज्या प्रत्येक चिन्हमालेला आहे, ती
$\Frm[L_0]$ मध्ये आहे, हे आपले विगमन गृहीतक आहे. रचनाक्रमिकेच्या
व्याख्येनुसार एकतर $\varphi_n$ आण्विक आहे किंवा पुढीलपैकी एखादी अट पूर्ण
करणारे $j,k < n$ अस्तित्वात असले पाहिजेत:
\begin{enumerate}
    \item $\varphi_n \ident \lnot \varphi_j$.%
    \item $\varphi_n \ident (\varphi_j \land \varphi_k)$.%
    \item $\varphi_n \ident (\varphi_j \lor \varphi_k)$.%
    \item $\varphi_n \ident (\varphi_j \lif \varphi_k)$.%
    \item $\varphi_n \ident (\varphi_j \liff \varphi_k)$.%
\end{enumerate}
आता प्रत्येक शक्यतेचा स्वतंत्र विचार करू. $\varphi_n$ आण्विक असेल, तर
$\varphi_n \in \Frm[L_0]$. त्याऐवजी $\varphi \ident
(\varphi_j \land \varphi_k)$ असे समजा.
\begin{quote}\small\textbf{स्रोतदुरुस्ती.} MRPL-002: गोठवलेल्या इंग्रजीत या एकाच शक्यतेत विन्यासात्मक अनन्यता-चिन्हाऐवजी पर्यायी सममूल्यता-चिन्ह आले आहे. आधीची स्पष्ट व्याख्या आणि इतर सर्व रचना-बाबी यांच्याशी जुळण्यासाठी येथे विन्यासात्मक अनन्यता वापरली आहे.\end{quote}
\ref{pl:syn:fseq:lem:fseq-init} नुसार
$\tuple{\varphi_0,\dotsc,\varphi_j}$ आणि $\tuple{\varphi_0,\dotsc,\varphi_k}$ या अनुक्रमे
$\varphi_j$ आणि~$\varphi_k$ साठी रचनाक्रमिका आहेत. या $\varphi$ च्या
रचनाक्रमिकेच्या उचित आरंभीच्या उपक्रमिका असल्यामुळे दोन्हींची लांबी
$n$ पेक्षा कमी आहे. म्हणून विगमन गृहीतकानुसार $\varphi_j$ आणि~$\varphi_k$ ही
$\Frm[L_0]$ मध्ये आहेत; त्यामुळे सूत्राच्या व्याख्येनुसार
$(\varphi_j \land \varphi_k)$ हेही त्यात आहे. इतर शक्यता समांतर युक्तिवादाने
सिद्ध होतात.
\end{proof}









\section{सत्य-मूल्यांकन आणि पूर्ति}\label{pl:syn:val:sec}

\begin{defn}[सत्य-मूल्यांकने]
``सत्य'' आणि ``असत्य'' ही दोन सत्यतामूल्ये असलेला संच
$\{\True, \False\}$ असू द्या. $\Lang{L_0}$ साठीचे
\emph{सत्य-मूल्यांकन} म्हणजे भाषेतील विधानीय चलांना
$\True$ किंवा $\False$ यांपैकी एक मूल्य देणारे फलन~$\pAssign{v}$;
म्हणजेच $\pAssign{v} \colon \PVar \to \{\True, \False \}$.
\end{defn}

\begin{defn}
  दिलेले सत्य-मूल्यांकन~$\pAssign{v}$ असताना, मूल्यन फलन
  $\pValue{v} \colon \Frm[L_0] \to \{\True, \False \}$ पुढीलप्रमाणे
  विगमनाने परिभाषित करू:
  \begin{align*}
    \pValue{v}(\lfalse) & = \False; \\
    \pValue{v}(\ltrue) & = \True; \\
    \pValue{v}(\Obj p_n) & = \pAssign{v}(\Obj p_n); \\
    \pValue{v}(\lnot \varphi) & = \begin{cases}
        \True & \text{जर } \pValue{v}(\varphi) = \False;\\
        \False & \text{अन्यथा.}
      \end{cases} \\
    \pValue{v}(\varphi \land \psi) & = \begin{cases}
        \True &
        \text{जर $\pValue{v}(\varphi) = \True$ आणि $\pValue{v}(\psi) = \True$;}\\
        \False &
        \text{जर $\pValue{v}(\varphi) = \False$ किंवा $\pValue{v}(\psi) = \False$}.
      \end{cases}\\
    \pValue{v}(\varphi \lor \psi) & = \begin{cases}
        \True &
        \text{जर $\pValue{v}(\varphi) = \True$ किंवा $\pValue{v}(\psi) = \True$;}\\
        \False &
        \text{जर $\pValue{v}(\varphi) = \False$ आणि $\pValue{v}(\psi) = \False$}.
      \end{cases}\\
    \pValue{v}(\varphi \lif \psi) & = \begin{cases}
        \True &
        \text{जर $\pValue{v}(\varphi) = \False$ किंवा $\pValue{v}(\psi) = \True$;}\\
        \False &
        \text{जर $\pValue{v}(\varphi) = \True$ आणि $\pValue{v}(\psi) = \False$}.
      \end{cases}\\
    \pValue{v}(\varphi \liff \psi) & = \begin{cases}
        \True &
        \text{जर $\pValue{v}(\varphi) = \pValue{v}(\psi)$;}\\
        \False &
        \text{जर $\pValue{v}(\varphi) \neq \pValue{v}(\psi)$}.
      \end{cases}
\end{align*}
\end{defn}

\begin{explain}
या बाबींना पुढील सत्यताकोष्टके अनुरूप आहेत:
\begin{center}

  \begin{tabular}{|c||c|} \hline
    $\varphi$ & $ \lnot \varphi$ \\
    \hline \hline
    $\True$ & $\False$ \\
    $\False$ & $\True$ \\
    \hline
  \end{tabular}


  \begin{tabular}{|cc||c|} \hline
    $\varphi$ & $\psi$ & $\varphi \land \psi$ \\
    \hline \hline
    $\True$ & $\True$ & $\True$ \\
    $\True$ & $\False$ & $\False$ \\
    $\False$ & $\True$ & $\False$ \\
    $\False$ & $\False$ & $\False$ \\
    \hline
  \end{tabular}


  \begin{tabular}{|cc||c|} \hline
    $\varphi$ & $\psi$ & $\varphi \lor \psi$ \\
    \hline \hline
    $\True$ & $\True$ & $\True$ \\
    $\True$ & $\False$ & $\True$ \\
    $\False$ & $\True$ & $\True$ \\
    $\False$ & $\False$ & $\False$ \\
    \hline
  \end{tabular}
%
\\[1em]
  \begin{tabular}{|cc||c|} \hline
    $\varphi$ & $\psi$ & $\varphi \lif \psi$ \\
    \hline \hline
    $\True$ & $\True$ & $\True$ \\
    $\True$ & $\False$ & $\False$ \\
    $\False$ & $\True$ & $\True$ \\
    $\False$ & $\False$ & $\True$ \\
    \hline
  \end{tabular}
  \begin{tabular}{|cc||c|} \hline
    $\varphi$ & $\psi$ & $\varphi \liff \psi$ \\
    \hline \hline
    $\True$ & $\True$ & $\True$ \\
    $\True$ & $\False$ & $\False$ \\
    $\False$ & $\True$ & $\False$ \\
    $\False$ & $\False$ & $\True$ \\
    \hline
  \end{tabular}
\end{center}
\end{explain}
\begin{prob}
$\Lang{L_0}$ मध्ये $\diamondsuit$ हा त्रिस्थानी संयोजक वाढवण्याचा विचार
करा; त्याचे मूल्यन पुढीलप्रमाणे दिले आहे:
\begin{gather*}
  \pValue{v}(\diamondsuit( \varphi , \psi , \chi ) ) = \begin{cases}
    \pValue{v}( \psi ) &
    \text{जर $\pValue{v}(\varphi) = \True$;}\\
    \pValue{v}( \chi ) &
    \text{जर $\pValue{v}(\varphi) = \False $}.
  \end{cases}
\end{gather*}
या संयोजकाचे सत्यताकोष्टक लिहा.
\end{prob}
\begin{thm}[स्थानिक निर्धारण]
\label{pl:syn:val:thm:LocalDetermination} $\pAssign{v_1}$ आणि
$\pAssign{v_2}$ ही अशी सत्य-मूल्यांकने आहेत की $\varphi$ मध्ये आस्थित असलेल्या
विधानीय चलांना ती समान मूल्ये देतात; म्हणजे, एखादे $\Obj p_n$ हे
सूत्र~$\varphi$ मध्ये आस्थित असेल तेव्हा $\pAssign{v_1}(\Obj p_n) =
\pAssign{v_2}(\Obj p_n)$. मग $\pValue{v_1}$ आणि $\pValue{v_2}$ हीही
$\varphi$ वर एकमत असतात; म्हणजे $\pValue{v_1}(\varphi) = \pValue{v_2}(\varphi)$.
\end{thm}
\begin{proof}
$\varphi$ वरील विगमनाने.
\end{proof}
\begin{defn}[पूर्ति]
\label{pl:syn:val:defn:satisfaction} आपण
  \emph{सत्य-मूल्यांकन~$\pAssign{v}$ कडून सूत्र~$\varphi$ ची पूर्ति},
  $\pSat{v}{\varphi}$, ही संकल्पना पुढीलप्रमाणे विगमनाने परिभाषित करू शकतो.
  (``$\pSat{v}{\varphi}$ नाही'' हे दर्शवण्यासाठी आपण $\pSat/{v}{\varphi}$ लिहितो.)
\begin{enumerate}
\item %
  \indcase{\varphi}{\lfalse}{$\pSat/{v}{\indfrm}$.}
\item %
  \indcase{\varphi}{\ltrue}{$\pSat{v}{\indfrm}$.}
\item \indcase{\varphi}{\Obj p_i}{$\pSat{v}{\indfrm}$ तेव्हा आणि केवळ
  तेव्हाच, जेव्हा $\pAssign{v}(\Obj p_i) = \True$.}
\item %
  \indcase{\varphi}{\lnot \psi}{$\pSat{v}{\indfrm}$ तेव्हा आणि केवळ
    तेव्हाच, जेव्हा $\pSat/{v}{\psi}$.}
\item %
  \indcase{\varphi}{(\psi \land \chi)}{$\pSat{v}{\indfrm}$ तेव्हा आणि केवळ
    तेव्हाच, जेव्हा $\pSat{v}{\psi}$ आणि $\pSat{v}{\chi}$.}
\item %
  \indcase{\varphi}{(\psi \lor \chi)}{$\pSat{v}{\indfrm}$ तेव्हा आणि केवळ
    तेव्हाच, जेव्हा $\pSat{v}{\psi}$ किंवा $\pSat{v}{\chi}$ (किंवा दोन्ही).}
\item %
  \indcase{\varphi}{(\psi \lif \chi)}{$\pSat{v}{\indfrm}$ तेव्हा आणि केवळ
    तेव्हाच, जेव्हा $\pSat/{v}{\psi}$ किंवा $\pSat{v}{\chi}$ (किंवा दोन्ही).}
\item %
  \indcase{\varphi}{(\psi \liff \chi)}{$\pSat{v}{\indfrm}$ तेव्हा आणि केवळ
    तेव्हाच, जेव्हा $\pSat{v}{\psi}$ आणि $\pSat{v}{\chi}$ ही दोन्ही, किंवा
    $\pSat{v}{\psi}$ आणि $\pSat{v}{\chi}$ यांपैकी कोणतेही नाही.}
\end{enumerate}
$\Gamma$ हा सूत्रांचा संच असेल, तर $\pSat{v}{\Gamma}$ तेव्हा आणि
केवळ तेव्हाच, जेव्हा प्रत्येक~$\varphi \in \Gamma$ साठी $\pSat{v}{\varphi}$.
\end{defn}
\begin{prop}\label{pl:syn:val:prop:sat-value}
  $\pSat{v}{\varphi}$ तेव्हा आणि केवळ तेव्हाच, जेव्हा
  $\pValue{v}(\varphi) = \True$.
\end{prop}
\begin{proof}
  $\varphi$ वरील विगमनाने.
\end{proof}
\begin{prob}
  \ref{pl:syn:val:prop:sat-value} सिद्ध करा.
\end{prob}
\section{चिन्हार्थविषयक संकल्पना}\label{pl:syn:sem:sec}
पुढील चिन्हार्थविषयक संकल्पना परिभाषित करू:
\begin{defn}
\begin{enumerate}
\item सूत्र~$\varphi$ हे \emph{पूर्ततायोग्य} असते, जर एखाद्या
  $\pAssign{v}$ साठी $\pSat{v}{\varphi}$; आणि कोणत्याही $\pAssign{v}$ साठी
  $\pSat{v}{\varphi}$ नसल्यास ते \emph{अपूर्ततायोग्य} असते;
\item सर्व सत्य-मूल्यांकने~$\pAssign{v}$ साठी $\pSat{v}{\varphi}$ असेल, तर
  सूत्र~$\varphi$ हे \emph{सर्वतः सत्य सूत्र} असते;
\item सूत्र~$\varphi$ हे पूर्ततायोग्य असेल, पण सर्वतः सत्य सूत्र नसेल,
  तर ते \emph{परिस्थितीवश} असते;
\item $\Gamma$ हा सूत्रांचा संच असेल, तर $\Gamma \Entails \varphi$
  (``$\Gamma$ पासून $\varphi$ तार्किकरीत्या निष्पन्न होते'') तेव्हा आणि केवळ
  तेव्हाच, जेव्हा $\pSat{v}{\Gamma}$ असलेल्या प्रत्येक
  सत्य-मूल्यांकन~$\pAssign{v}$ साठी $\pSat{v}{\varphi}$.
\item $\Gamma$ हा सूत्रांचा संच असेल, तर एखाद्या
  सत्य-मूल्यांकन~$\pAssign{v}$ साठी $\pSat{v}{\Gamma}$ असल्यास $\Gamma$
  हा \emph{पूर्ततायोग्य} असतो; अन्यथा $\Gamma$ हा
  \emph{अपूर्ततायोग्य} असतो.
\end{enumerate}
\end{defn}
\begin{prob}
 पुढील चार सूत्रांपैकी प्रत्येक सूत्र (a)~पूर्ततायोग्य,
 (b)~सर्वतः सत्य सूत्र आणि (c)~परिस्थितीवश आहे की नाही ते ठरवा.
\begin{enumerate}
  \item \( ( \Obj p_0 \lif ( \lnot \Obj p_1 \lif \lnot \Obj p_0 ) ) \).
  \item \( ( ( \Obj p_0 \land \lnot \Obj p_1 ) \lif ( \lnot \Obj p_0 \land \Obj p_2 )) \liff ( ( \Obj p_2 \lif \Obj p_0 ) \lif ( \Obj p_0 \lif \Obj p_1 )) \).
  \item \( ( \Obj p_0 \liff \Obj p_1 ) \lif ( \Obj p_2 \liff \lnot \Obj p_1 ) \).
  \item \( (( \Obj p_0 \liff ( \lnot \Obj p_1 \land \Obj p_2 )) \lor ( \Obj p_2 \lif ( \Obj p_0 \liff \Obj p_1 ))) \).
\end{enumerate}
\end{prob}
\begin{prop}
\label{pl:syn:sem:prop:semanticalfacts}
\begin{enumerate}
\item $\varphi$ हे सर्वतः सत्य सूत्र तेव्हा आणि केवळ तेव्हाच असते, जेव्हा
  $\emptyset \Entails \varphi$;
\item $\Gamma \Entails \varphi$ आणि $\Gamma \Entails \varphi \lif \psi$ असल्यास
  $\Gamma \Entails \psi$;
\item $\Gamma$ पूर्ततायोग्य असेल, तर $\Gamma$ चा प्रत्येक सांत उपसंचही
  पूर्ततायोग्य असतो;
\item \label{pl:syn:sem:def:monotonicity} एकस्वनिकता: $\Gamma \subseteq \Delta$
  आणि $\Gamma \Entails \varphi$ असल्यास $\Delta \Entails \varphi$ हेही असते;
\item \label{pl:syn:sem:def:Cut} संक्रमकता: $\Gamma \Entails \varphi$ आणि
  $\Delta \cup \{ \varphi\} \Entails \psi$ असल्यास $\Gamma \cup \Delta \Entails
  \psi$.
\end{enumerate}
\end{prop}
\begin{proof}
सरावासाठी.
\end{proof}
\begin{prob}
\ref{pl:syn:sem:prop:semanticalfacts} सिद्ध करा.
\end{prob}
\begin{prop}\label{pl:syn:sem:prop:entails-unsat}
  $\Gamma \Entails \varphi$ तेव्हा आणि केवळ तेव्हाच, जेव्हा
  $\Gamma \cup \{\lnot \varphi\}$ अपूर्ततायोग्य असतो.
\end{prop}
\begin{proof}
सरावासाठी.
\end{proof}
\begin{prob}
\ref{pl:syn:sem:prop:entails-unsat} सिद्ध करा.
\end{prob}
\begin{thm}[चिन्हार्थविषयक निगमन प्रमेय]
  \label{pl:syn:sem:thm:sem-deduction} $\Gamma \Entails \varphi \lif \psi$ तेव्हा आणि केवळ
  तेव्हाच, जेव्हा $\Gamma \cup \{\varphi\} \Entails \psi$.
\end{thm}
\begin{proof}
सरावासाठी.
\end{proof}
\begin{prob}
\ref{pl:syn:sem:thm:sem-deduction} सिद्ध करा.
\end{prob}
